\section{DSPy 复杂优化与约束编码}

最后一部分讲述 DSPy 框架如何处理组合优化里难以表达的非线性约束，如 nanophotonics 的 birefringent 设计。

\subsection{潜在线性化 \& Green Descent}

对于复杂约束，团队先预测一个 cost vector \(C\)，再把原问题变成线性子问题，最后用 solver 生成 \(X^\star\)。整个过程构成一张计算图，从 description \(Y\) 到 \(C(Y)\)、再到 solver 输出，loss 是与原时域目标之间的差距，通过 gradient descent 反传训练 C 的表示。
\begin{figure}[H]
\centering
\includegraphics[width=\textwidth]{frame-latent-opt.jpg}
\caption{DSPy 中从描述 \(Y\) 推导 \(C\)，再调用 solver 产出 \(X^\star\) 的计算图。 \protect\footnotemark}
\end{figure}
\footnotetext{视频画面时间区间：00:45:05--00:45:25。}
\begin{knowledgebox}{Latent Optimization Pipeline}
1）从描述预测 latent cost \(C\)；2）固定 \(C\) 后用可微 solver 近似求解；3）将输出回插至物理 loss；4）用梯度下降训练 \(C\)。这种闪电般的反解比直接求解原始非线性问题更稳定。
\end{knowledgebox}
\begin{importantbox}{Green Descent 的分层步骤}
先用轻量级 solver 解析线性化子问题，生成 \(X^\star\)，然后把输出过 backprop 回 latent cost，再用梯度下降调整 \(C\)。这种两层迭代（solver + gradient update）在数值上避免了直接在原始非线性空间里爬山的震荡。
\end{importantbox}

\subsection{光学/制造约束}

在 80\u00d780 网格的光学设计里，每个格点只能写 0/1，而且制造工艺要求不能出现孤立的点、必须满足笔刷大小。还要考虑频率响应、干涉路径等多维度约束，因此把问题分解到多个 benchmark（beam splitter、wavelength multiplexer）上验证。
\begin{warningbox}{制造约束隐藏陷阱}
约束不仅有代数形式，还包括工程可制造性：刷具大小导致的线性约束、不可行的孤立像素、以及发光路径的拓扑限制。任何忽略这些的优化都会在制造阶段报废。
\end{warningbox}

\subsection{本章小结}

通过 latent linearization + solver 回写 DSPy 让复杂约束的学习从 ``无穷搜索空间'' 降维到 ``可微参数''，同时把制造约束纳入评价体系。

